\small
\setlength{\tabcolsep}{4pt}
\renewcommand{\arraystretch}{1.2}

\begin{longtable}{
    >{\RaggedRight\arraybackslash}p{3.0cm}
    >{\RaggedRight\arraybackslash}p{4.2cm}
    >{\RaggedRight\arraybackslash}p{5.0cm}
    >{\RaggedRight\arraybackslash}p{1.5cm}
}
\toprule
\textbf{Panel} & \textbf{Revision Objective} & \textbf{Action Taken} & \textbf{Page(s)} \\
\midrule
\endfirsthead

\multicolumn{4}{c}{{Continued from previous page}} \\
\toprule
\textbf{Panel} & \textbf{Revision Objective} & \textbf{Action Taken} & \textbf{Page(s)} \\
\midrule
\endhead

\midrule
\multicolumn{4}{r}{Continued on next page} \\
\endfoot

\bottomrule
\endlastfoot

% --- Asst. Prof. Julieto Perez (4 rows) ---
\multirow{4}{*}{\parbox[t]{2.8cm}{\raggedright Asst. Prof. Julieto Perez, MSc}}
& Reframe thesis as operational/optimization study, not full ecological/behavioral realism.
& Added explicit disclaimers in Abstract, Sections 1.1, 1.5, and 3.1 stating that the ABM is an abstract testbed for policy optimization.
& 8, 24–27, 29–31, 46–47 \\
\cline{2-4}
& Add ablation study comparing baseline vs. PPO and other heuristics.
& Created Section 4.3.3 with Table 4.6 comparing Random, Dirichlet, Greedy, Vanilla DRL, and HuDRL; included statistical tests (t‑test, Cohen’s d).
& 102–104 \\
\cline{2-4}
& Define clear performance metrics and conduct sensitivity analysis.
& Defined four metrics (Section 3.5.5); performed Global Sensitivity Analysis (Sobol method, Section 4.2) and tipping‑point robustness test (Section 4.2.3).
& 86, 95–99 \\
\cline{2-4}
& Adjust tone and limitations: acknowledge psychometric simplifications and lack of predictive validity.
& Added extended limitations in Sections 1.5 and 5.3; consistently used qualified language (“within the simulation,” “conditional on assumptions”).
& 29–31, 108 \\
\midrule

% --- Asst. Prof. Liezil Daberao (3 rows) ---
\multirow{3}{*}{\parbox[t]{2.8cm}{\raggedright Asst. Prof. Liezil Daberao, MSc}}
& Explain source, logic, and construction of Table 4.4 (Sobol indices).
& Added pre‑table paragraph in Section 4.2 detailing parameter selection criteria, literature sources (Paigalan et al., Meng et al.), and interview grounding.
& 95–96 \\
\cline{2-4}
& Correct and interpret Table 4.7; ensure numbers align with simulation results.
& Verified and corrected compliance numbers; added interpretative text after the table explaining each regime’s performance (Status Quo, Pure Enforcement, Pure Incentives, HuDRL).
& 113–114 \\
\cline{2-4}
& Address numeric indices that exceed theoretical upper bound (e.g., $S_T = 1.06$).
& Explicitly acknowledged finite‑sample estimation artifact in Section 4.2; explained that values slightly above 1.0 indicate dominant interaction effects.
& 97 \\
\midrule

% --- Asst. Prof. Jennifer Joyce Montemayor (4 rows) ---
\multirow{4}{*}{\parbox[t]{2.8cm}{\raggedright Asst. Prof. Jennifer Joyce Montemayor, MSc}}
& Write full reward function as a mathematical equation with all multipliers.
& Presented Equation (3.21) in Section 3.4.4; provided component‑by‑component explanation.
& 73–77 \\
\cline{2-4}
& Create a parameter table mapping each weight to the reward function.
& Added Table 3.3 (Parameter Mapping) with ranges, values, and effect descriptions.
& 77 \\
\cline{2-4}
& For each reward component, describe behavior encouraged/discouraged and acknowledge subjectivity.
& Extended Section 3.4.4 with explicit statements for each term; identified the waste penalty and attitude bonus as the most subjective components.
& 73–77 \\
\cline{2-4}
& Address sensitivity to thresholds (e.g., 0.5) and local minima.
& Added subsection “Sensitivity and Local Minima Concerns” within Section 3.4.4; tested neutral point values (0.3, 0.7) and entropy coefficients (Table 3.2).
& 75–76 \\
\midrule

% --- Asst. Prof. John Gieveson Iglupas (3 rows) ---
\multirow{3}{*}{\parbox[t]{2.8cm}{\raggedright Asst. Prof. John Gieveson Iglupas, MPA}}
& Frame problem as policy implementation/governance issue; discuss street‑level bureaucracy, political dynamics, CSOs.
& Added discussion in Section 4.6.1 (“Computational optimality versus political reality”) and Section 5.2 (Policy Brief), citing Carpio (2025), IBON (2026), PIDS (2024, 2025), and acknowledging CSO monitoring role.
& 98–102, 126–129 \\
\cline{2-4}
& Clarify core problem (low compliance, budget constraints, enforcement gaps) and link model insights to policy levers.
& Created Section 4.6 (Policy Implications) specifying tripartite failure and policy levers identified by the model (sequential saturation, pulsed enforcement, catalytic incentives).
& 98–102 \\
\cline{2-4}
& Distinguish empirical facts from simulated assumptions; emphasize interdisciplinary nature.
& Added explicit distinction in Section 4.2 (Sobol indices as model properties) and Section 4.5 (conditional findings); discussed interdisciplinary framing in Section 5.3.
& 97, 117–118, 108 \\

\end{longtable}